% =============================================================
%  OpenPhysicist Project Cover --- public template (for \input reuse)
%  Depends: fontspec, xeCJK, xcolor, tikz, fontawesome5
%           \usetikzlibrary{positioning,calc}
% =============================================================

\definecolor{coverprimary}{HTML}{0D1B2A}
\definecolor{coveraccent}{HTML}{C9A227}
\definecolor{coverteal}{HTML}{0E7C7B}
\definecolor{coverwarm}{HTML}{B85C38}
\definecolor{coverdark}{HTML}{1F2937}
\definecolor{covermuted}{HTML}{64748B}
\definecolor{goldpanel}{RGB}{255,248,230}
\definecolor{badgequantum}{HTML}{0E7C7B}
\definecolor{badgesym}{HTML}{4C5FD5}
\definecolor{badgenuclear}{HTML}{B85C38}
\definecolor{badgephil}{HTML}{8D6E63}

\newcommand{\openphysicistslide}{%
\begin{frame}[plain]
\begin{tikzpicture}[remember picture, overlay]
  \fill[coverprimary!4]   (current page.north west) rectangle (current page.south east);
  \fill[coverdark, opacity=0.06] (current page.south west) rectangle ([yshift=0.4cm]current page.south east);
  \draw[coverprimary, opacity=0.30, line width=0.8pt] ([yshift=0.4cm]current page.south west) -- ([yshift=0.4cm]current page.south east);

  \node[anchor=center, font=\fontsize{24}{30}\selectfont\bfseries, text=coverprimary!90!black]
    at ([yshift=3.9cm]current page.center) {OpenPhysicist 已开放源码};
  \node[anchor=center, font=\fontsize{14}{18}\selectfont, text=coverprimary!60!black]
    at ([yshift=3.05cm]current page.center) {为 2,000+ 位物理学家立传};

  \draw[coveraccent!45, line width=1.0pt] ([yshift=2.45cm, xshift=-5.0cm]current page.center)
    -- ([yshift=2.45cm, xshift=5.0cm]current page.center);

  \node[anchor=center, font=\fontsize{8.2}{11.5}\selectfont, text=coverdark!72, align=center, text width=13cm]
    at ([yshift=1.55cm]current page.center)
    {从牛顿、麦克斯韦、爱因斯坦，到费曼、杨振宁、维格纳，\\[2pt]
     OpenPhysicist 正在为物理学家建立开放、持续演进的物理学人物档案。};

  \node[anchor=center, font=\fontsize{8}{11}\selectfont\bfseries, text=coverprimary!82!black, align=center]
    at ([yshift=0.45cm]current page.center)
    {这不只是人物名单，而是一部由社区共同构建的现代物理学人物史：};

  \node[anchor=north west, font=\fontsize{7}{9.2}\selectfont, text=coverdark!78, align=left]
    at ([xshift=-6.3cm, yshift=-0.05cm]current page.center)
    {\textcolor{badgequantum}{\faBook}\enspace\textbf{2,000+ 物理学家} --- 持续收录、整理与补充\\[4pt]
     \textcolor{coveraccent}{\faMap}\enspace\textbf{人物生平} --- 教育经历、学术生涯、重要节点\\[4pt]
     \textcolor{badgesym}{\faTh}\enspace\textbf{物理贡献} --- 定理、理论、方法与历史影响};
  \node[anchor=north west, font=\fontsize{7}{9.2}\selectfont, text=coverdark!78, align=left]
    at ([xshift=0.3cm, yshift=-0.05cm]current page.center)
    {\textcolor{badgenuclear}{\faTrophy}\enspace\textbf{荣誉与传承} --- 诺贝尔奖、费米奖、普朗克奖章等\\[4pt]
     \textcolor{coverteal}{\faGlobe}\enspace\textbf{物理谱系} --- 学派、师承、合作与思想传承\\[4pt]
     \textcolor{badgephil}{\faFlask}\enspace\textbf{开放共建} --- 欢迎任何人补充、纠错、完善};

  \node[draw=coveraccent!50, fill=goldpanel, rounded corners=5pt, inner sep=9pt, text width=12.0cm, align=center]
    at ([yshift=-2.35cm]current page.center)
    {{\fontsize{8.6}{11}\selectfont\bfseries\color{coverdark!88} 物理学史，不应该只存在于少数人的书架上。}\\[2pt]
     {\fontsize{7.5}{9.5}\selectfont\color{covermuted} 让每一位物理学家的故事都可以被发现、被理解、被连接。}};

  \node[anchor=south, font=\fontsize{7.8}{9.8}\selectfont\bfseries, text=coveraccent!80!black]
    at ([yshift=0.64cm]current page.south) {\faIcon{star}\enspace Star OpenPhysicist · 加入共建};
  \node[anchor=south, font=\fontsize{6.2}{7.8}\selectfont, text=covermuted]
    at ([yshift=0.38cm]current page.south) {欢迎 Star、Issue、PR，共同建设一部开放的物理学家人物史。};
\end{tikzpicture}
\end{frame}
}
